\documentclass[11pt,a4paper]{article}
\usepackage{fontspec,amsmath,amssymb,amsthm,microtype}
\setmainfont{lmroman10-regular.otf}[BoldFont=lmroman10-bold.otf,ItalicFont=lmroman10-italic.otf,BoldItalicFont=lmroman10-bolditalic.otf]
\usepackage[margin=25mm]{geometry}
\usepackage[unicode,hidelinks]{hyperref}
\hypersetup{pdfauthor={LangevinLean Project}}
\setlength{\emergencystretch}{3em}
\setlength{\parskip}{0.4em}
\setlength{\parindent}{0pt}
\allowdisplaybreaks[2]
\raggedbottom
\title{\texorpdfstring{Uniform Moreau gaps in scaled particle coordinates}{Uniform Moreau gaps in scaled particle coordinates}}
\author{LangevinLean Project}
\date{}
\begin{document}\maketitle
Draft mathematical exposition; these notes may contain mistakes.


The lifted envelope gradient has norm at most \(NK\), and its uniform potential gap is at most \(N\eta K^2/2\). These follow by summing the single-pair value estimates and are different from gradient variation with constant involving \(1/\eta\). Write \(F_g(z)=\sum_i g(J_i z)\) with \(N>0\). If \(g\) is convex, bounded below, and \(K\)-Lipschitz, the Moreau value Lipschitz bound yields a gradient-magnitude estimate uniform over positive smoothing parameters.


\section{\texorpdfstring{The lifted value bound and envelope-\allowbreak{}gradient magnitude}{The lifted value bound and envelope-gradient magnitude}}


The averaged potential has value Lipschitz constant \(K\). Multiplication by \(N\) gives \(\operatorname{Lip}(F_g)\le NK\). Since the envelope retains single-pair value constant \(K\), its lifted sum \(F_{g^\eta}\) has constant \(NK\). Differentiability therefore bounds its gradient norm by \(NK\) for every \(\eta>0\).



Equivalently multiplication by \(N\) gives lifted value constant \(NK\). The envelope’s single-pair constant \(K\) yields \(\|\nabla V_{N,\eta}^{(2)}\|\le NK\). This uses the norm bound for the derivative of a Lipschitz function, rather than the gradient Lipschitz constant \(N/\eta\), and is consequently uniform as \(\eta\downarrow0\). In the averaged notation the constants are \(K\) and \(NK\) for the full sum, giving \(\|\nabla F_{g_2^\eta}\|\le NK\) independently of \(\eta\). Rapid spatial variation of the gradient does not change this bound on its magnitude.




Earlier results used here: \cite[Section~3.1.1]{LangevinLean2026IPLA}.

\section{\texorpdfstring{The summed envelope gap and split-\allowbreak{}potential identity}{The summed envelope gap and split-potential identity}}


Apply \(0\le g-g^\eta\le\eta K^2/2\) at each pair and sum the finite inequalities:

\[0\le F_g(z)-F_{g^\eta}(z)\le N\eta K^2/2.\]

Distributivity also gives the function identity \(F_{g_1+g_2^\eta}=F_{g_1}+F_{g_2^\eta}\). In nonsmooth-component notation the bound \(0\le g_2-g_2^\eta\le\eta K^2/2\) has \(N\) summands, so

\[0\le V_N^{(2)}-V_{N,\eta}^{(2)}\le N\eta K^2/2.\]

Expanding \(V_N(g_1+g_2^\eta)\) and distributing the sum gives \(V_N(g_1)+V_N(g_2^\eta)\). Each nonnegative gap is at most \(\eta K^2/2\); summing \(N\) such gaps proves the displayed total gap directly. Differentiating the sum identity, under the established regularity, separates the smooth and envelope reference forces in the splitting argument.




Earlier results used here: \cite[Section~3.1.1]{LangevinLean2026IPLA}.

\begin{thebibliography}{ULA-LSI}
\bibitem[IPLA-SC]{LangevinLean2026IPLA}
LangevinLean Project. \emph{Interacting Particle Langevin under Strong Convexity}. LL-03, 2026. Draft report, content revision c8b5d67ceb43. \url{https://odakyildiz.github.io/LangevinLean/docs/pdf/IPLA/FullProof.pdf}.
\end{thebibliography}

\end{document}
